\question{}

پرسش‌های درست و نادرست:

\begin{ابجد}
\item \textbf{درست}؛ سیاست استخراج‌شده تنها به مقادیر تقریبی ارزش حالت‌ها و ترتیب آنها بستگی دارد و نه مقادیر دقیق. بنابراین معمولا پیش از اینکه ارزش‌ها به همگرایی دقیق برسند، سیاست بهینه به دست می‌آید.

\item \textbf{نادرست}؛ الگوریتم \lr{Q-Learning} همواره به مقادیر بهینه همگرا می‌شود، حتی اگر بر اساس سیاست بهینه عمل نکنیم.

\item \textbf{درست}؛ مدت زمان اجرای هر \lr{iteration} از الگوریتم \lr{Value Iteration}، از مرتبه‌ی $O(|S|^2|A|)$، و برای الگوریتم \lr{Policy Iteration}، $O(|S|^3 + |S|^2|A|)$ می‌باشد. بنابراین اگر $|S| \approx |A|$، پیچیدگی زمانی هر \lr{iteration} از دو الگوریتم از یک مرتبه‌ی یکسان می‌باشد.

\item \textbf{درست} (در صورتی که مسئله دارای افق زمانی نامتناهی و $\gamma < 1$ باشد)؛ فرض می‌کنیم پاداش‌های جدید $R'(s,a,s') = R(s,a,s') + c$ باشند. طبق معادلات بلمن می‌دانیم که ارزش‌های بهینه به این صورت تعریف می‌شوند:
\begin{align*}
    V^*(s) = \max_a \sum_{s'} T(s, a, s') \left[ R(s, a, s') + \gamma V^*(s') \right]
\end{align*}
و برای پاداش‌های جدید داریم:
\begin{align*}
    V_{\text{new}}^*(s) = \max_a \sum_{s'} T(s, a, s') \left[ R(s, a, s') + c + \gamma V_{\text{new}}^*(s') \right]
\end{align*}
حدس می‌زنیم که ارزش تمام حالت‌ها به یک میزان ثابت افزایش می‌یابد ($V_{\text{new}}^*(s) = V^*(s) + \Delta$). در این صورت داریم:
\begin{align*}
     & V_{\text{new}}^*(s) = \max_a \sum_{s'} T(s, a, s') \left[ \left( R(s, a, s') + \gamma V^*(s') \right) + (c + \gamma \Delta) \right]          \\
     & V_{\text{new}}^*(s) = \max_a \left( \sum_{s'} T(s, a, s') [R(s, a, s') + \gamma V^*(s')] + \sum_{s'} T(s, a, s') (c + \gamma \Delta) \right) \\
     & V_{\text{new}}^*(s) = \max_a \left( \sum_{s'} T(s, a, s') [R(s, a, s') + \gamma V^*(s')] + (c + \gamma \Delta) \right)                       \\
     & V_{\text{new}}^*(s) = V^*(s) + c + \gamma \Delta
\end{align*}
فرض کرده بودیم که این افزایش به میزان $\Delta$ باشد. اکنون داریم:
\begin{align*}
    \Delta = c + \gamma \Delta \implies \Delta = \frac{c}{1 - \gamma}
\end{align*}
از آنجا که طبق رابطه‌ی $\pi^*_{\text{new}}(s) = \arg\max_a \sum_{s'} T(s, a, s') \left[ R'(s, a, s') + \gamma V_{\text{new}}^*(s') \right]$ برای پیدا کردن سیاست بهینه باید بین مقادیر مختلف بیشینه بگیریم، با اضافه شدن یک مقدار ثابت به تمام ارزش‌ها، سیاست بهینه تغییر نمی‌کند.

\item \textbf{نادرست}؛ از نظر شهودی، ضریب تخفیف را تعریف می‌کنیم تا \lr{reward} های زودتر ارزش بیشتری داشته باشند. حال اگر دو \lr{MDP} یکسان با ضریب تخفیف‌های متفاوت داشته باشیم ($\gamma$ اولی بیشتر باشد)، الگوریتم دوم تلاش می‌کند مسیرهایی را انتخاب کند که در ابتدا \lr{reward} بیشتری دارند، ولی الگوریتم اول کمتر به این موضوع اهمیت می‌دهد. بنابراین ممکن است سیاست بهینه‌ی این دو با هم متفاوت باشد.

\item \textbf{درست}؛ در مسائل با افق محدود، تصمیم بهینه علاوه بر حالت فعلی، به اینکه چقدر زمان باقی مانده نیز بستگی دارد. به عنوان مثال اگر تنها یک گام باقی مانده باشد، باید پاداش همان مرحله را بیشینه کنیم، اما اگر مراحل بیشتری باقی مانده باشد، ممکن است اکشنی را انتخاب کنیم که پاداش فوری کمتری داشته باشد، اما کمک کند که در مراحل بعدی پاداش‌های بهتری بگیریم.

\end{ابجد}

پرسش‌های با پاسخ کوتاه:

\begin{ابجد}
\item در روش‌های \lr{on-policy}، سیاستی که در حال ارزیابی و بهبود آن هستیم، دقیقاً همان سیاستی است که از آن برای تعامل با محیط و جمع‌آوری داده‌های نمونه استفاده می‌شود. به عبارت دیگر، عامل از تجربیات حاصل از سیاست فعلی خود یاد می‌گیرد. اما در روش‌های \lr{off-policy}، سیاستی که ارزیابی و بهینه‌سازی می‌شود با سیاستی که برای جمع‌آوری داده‌ها استفاده شده متفاوت است؛ در این حالت داده‌ها می‌توانند توسط یک سیاست مجزا (که حتی می‌تواند کاملاً تصادفی یا صرفاً برای اکتشاف محیط باشد) تولید شده باشند.

\item در روش‌های \lr{model-based}، عامل ابتدا تلاش می‌کند یک مدل صریح و تجربی از رفتار محیط (شامل احتمال انتقال بین حالت‌ها و پاداش‌ها) را با روش‌های آماری یاد بگیرد و سپس با استفاده از روش‌های برنامه‌ریزی پویا مانند \lr{Value Iteration} یا \lr{Policy Iteration} آن را حل کند. در مقابل، در روش‌های \lr{model-free} عامل نیازی به یادگیری یا دانستن مدل محیط ندارد، بلکه مستقیماً از طریق تعامل، تجربه و نمونه‌برداری از مسیرها، ارزش حالت‌ها یا خود سیاست را به‌روزرسانی و بهینه‌سازی می‌کند.

\end{ابجد}